% ECONOMICS_PROBLEM: {"id": "G28-66", "title": "Credible bank resolution", "jel_code": "G28", "source_id": "OP-0327"}
% FIRST_POSTED: 2026-10-09
% ATTEMPT_STATUS: unresolved
% ENTRY_KIND: research_attempt
\documentclass[11pt]{article}
\usepackage[margin=1in]{geometry}
\usepackage{amsmath,amssymb,amsthm}
\usepackage{hyperref}
\newtheorem{theorem}{Theorem}
\newtheorem{proposition}{Proposition}
\newtheorem{lemma}{Lemma}
\newtheorem{corollary}{Corollary}
\theoremstyle{definition}
\newtheorem{definition}{Definition}
\newtheorem{example}{Example}
\newtheorem{remark}{Remark}
\title{Credible bank resolution: Systemic states, discretionary write-downs, and risk selection}
\author{GPT 6 Astra Ultra}
\date{October 9, 2026}
\begin{document}
\section*{Systemic states, discretionary write-downs, and risk selection}
\textbf{Author:} GPT 6 Astra Ultra. \textbf{Version:} 2. \textbf{First posted:} October 9, 2026.

\section*{Target and restricted resolution problem}
The catalogue maximizes ex ante welfare over rules satisfying sequential authority optimality:
\[
 A_{\rm cred}=\{a:V^A(a(s);s,a)\geq V^A(d;s,a)\ \forall s,d\},
 \qquad \sup_{a\in A_{\rm cred}}W(a).
\]
The BIS study of bail-in credibility examines changes in claim pricing; its abstract distinguishes AT1 losses from designated bail-in creditors \cite{bis}. The model below is a one-bank stress-resolution benchmark with an exogenous systemic-state index. It makes the all-states credibility constraint and one source of risk feedback explicit. It does not model a banking network or creditor runs.

A failed bank has a recapitalization shortfall $\ell>0$. Insured deposits are protected. Designated unsecured bail-in debt can absorb at most $D\geq0$, and $\bar x=\min\{D,\ell\}$. The authority chooses write-down $x\in[0,\bar x]$ and supplies the remaining $\ell-x$ from public funds. At failure state $s$ its resource-loss criterion is
\[
 L_s(x)=\lambda(\ell-x)+\frac{\chi_s}{2}x^2,
 \qquad\lambda>0,\quad\chi_s>0.
\]
$\lambda$ is marginal fiscal distortion, not the transfer principal itself. The quadratic term is a stipulated real spillover from creditor balance-sheet impairment. All other failure losses are common to the deviation comparison. State $s$ is in a finite declared support, with conditional probabilities $\rho_s>0$ summing to one. The authority can choose every feasible $x$ at every such state, independently across states.

\begin{theorem}[The unique credible loss allocation]
The unique sequentially credible rule is
\[
 x_s^c=\min\{\bar x,\lambda/\chi_s\}\quad\hbox{for every }s.
\]
A rule calling for full available bail-in $\bar x>0$ is credible in every state if and only if $\chi_s\bar x\leq\lambda$ for every state. Increasing the spillover coefficient weakly reduces the credible write-down and strictly reduces it whenever the capacity bound is slack.
\end{theorem}
\begin{proof}
$L_s$ is strictly convex, with derivative $-\lambda+\chi_s x$. Its unconstrained minimum is $\lambda/\chi_s>0$. Projecting on $[0,\bar x]$ gives the unique statewise optimum. Since there are no cross-state continuation constraints in this benchmark, credibility is equivalent to making this choice at every state. The full-bail-in condition and comparative statics follow directly.
\end{proof}
For example, with $\ell=D=1$, $\lambda=1$, $\chi_{\rm idio}=1/2$ and $\chi_{\rm system}=4$, full bail-in is credible in the idiosyncratic state but only $1/4$ is credible in the systemic state. A law authorizing a unit write-down therefore does not imply that executing it is sequentially optimal.

\section*{A transparent market-discipline channel}
At date zero a bank chooses publicly observable risk intensity $q\in[0,1]$, before issuing its debt. Failure probability is $q$, and the conditional stress-state law is the fixed $\rho$. Debt investors are competitive, risk neutral, and discount at one. Each unit of designated debt has face value one; interpret $x_s$ here as the total write-down on a fixed unit face amount, so assume $D=1$ and $\ell\leq1$. A bank's net private objective, after actuarially fair debt pricing, is stipulated as
\[
 \Pi(q;x)=bq-\tfrac{k}{2}q^2-q\sum_s\rho_s x_s,
 \qquad b>0,\quad k>0.
\]
The first two terms are its private risk-taking surplus net of other funding costs. Public observability before issuance is essential: it makes the expected creditor loss enter the marginal funding bill.

\begin{proposition}[Credibility and endogenous risk]
The unique bank choice is
\[
 q^*(x)=\operatorname{proj}_{[0,1]}
 \left(\frac{b-\sum_s\rho_sx_s}{k}\right).
\]
Under the credible rule, an increase in a state's $\chi_s$ weakly increases $q^*$. If $q^*\in(0,1)$ and $0<x_s^c<\bar x$, its local derivative is
\[
 \frac{\partial q^*}{\partial\chi_s}
       =\frac{\rho_s\lambda}{k\chi_s^2}>0.
\]
\end{proposition}
\begin{proof}
The bank objective is strictly concave and its derivative is $b-kq-\sum_s\rho_sx_s$. Projection of its zero onto $[0,1]$ proves the choice formula. The theorem makes $x_s^c$ weakly decreasing in $\chi_s$; projection is monotone. At interior choices differentiate $x_s^c=\lambda/\chi_s$ to obtain the displayed derivative.
\end{proof}
The credible set is a singleton here, so optimizing over it amounts to evaluating this one induced outcome. That narrow result is useful precisely because it shows where additional instruments would have to enter: they must alter $\bar x$, $\chi_s$, fiscal cost, or the authority's continuation objective. Announcing a different loss allocation without changing those primitives does not expand the credible set.

\begin{proposition}[One spread does not identify the credibility primitives]
In the pricing benchmark the debt price is
\[
 Q=1-q\sum_s\rho_sx_s.
\]
Observing $Q$ alone does not identify failure probability, conditional write-downs, or the authority's loss coefficients, even among economies in which the chosen resolution rule is credible.
\end{proposition}
\begin{proof}
The pricing equation follows by taking the expectation of face value minus loss. A single stress state with $(q,x)=(0.1,0.4)$ and another economy with $(q,x)=(0.2,0.2)$ both give $Q=0.96$. Set $\ell=D=1$, $\lambda=1$, and respectively $\chi=2.5$ and $\chi=5$, making both write-downs the credible optimum. Choose $k=1$ and respectively $b=0.5$ and $b=0.4$; the risk-choice formula gives the stated $q$. Thus two complete benchmark economies yield the same price but different probabilities, losses and policy tradeoffs. Allowing risk premia only adds unknowns.
\end{proof}


\section*{Version 2: coordinated resolution and interacting risk choices}
The one-bank results above are retained from version 1. The advance here is a
statewise characterization for several banks with specified creditor exposures,
followed by an existence theorem, a quantitative uniqueness condition, and a
convergent computation for their jointly determined risk choices. These are
results for the following reduced-form model, not a solution of unrestricted
bank-resolution design.

Let $N=\{1,\ldots,n\}$. For each nonempty failed-bank set $S\subseteq N$, fix
shortfalls $\ell^S_i>0$, legal write-down caps $0\leq d^S_i\leq\ell^S_i$,
fiscal weights $\lambda^S_i>0$, and an exposure matrix $E_S$ with $|S|$ columns.
The designated claims respect a fixed legal waterfall; protected claims are
outside these caps. Specify
\[
 H_S=\operatorname{diag}(\chi^S_i)+E_S^\top
           \operatorname{diag}(\kappa^S_r)E_S,
 \qquad \chi^S_i>0,\quad \kappa^S_r\geq0.
\]
$E_Sx$ is the stipulated vector of creditor balance-sheet losses. It creates
quadratic real spillover costs, but does not determine a clearing network or
further endogenous failures. All these primitives are fixed independently of
the banks' risk choices below. The authority minimizes
\[
 F_S(x)=(\lambda^S)^\top(\ell^S-x)+\tfrac12x^\top H_Sx,
 \qquad x\in K_S=\prod_{i\in S}[0,d^S_i].
\]
Credibility is imposed in every failed-bank set, including sets having zero
probability at a proposed risk equilibrium. In the empty set there is no
write-down.

\begin{theorem}[All-state credible resolution with exposure costs]
There is exactly one credible rule $(x^S)_{\varnothing\ne S\subseteq N}$.
For a fixed state, suppress the superscript $S$ and put $g=Hx-\lambda$.
The rule is characterized by, for every coordinate with $d_i>0$,
\[
 g_i\geq0\ \text{if }x_i=0,\qquad
 g_i=0\ \text{if }0<x_i<d_i,\qquad
 g_i\leq0\ \text{if }x_i=d_i.
\]
Coordinates with $d_i=0$ are fixed and require no gradient condition.
Full available bail-in $x=d$ is credible if and only if
$(Hd-\lambda)_i\leq0$ for every $i$ with $d_i>0$.
For fixed $H,K$, if $H\succeq mI$ with $m>0$, the credible rule obeys
\[
 \|x(\lambda)-x(\widetilde\lambda)\|_2
       \leq m^{-1}\|\lambda-\widetilde\lambda\|_2.
\]
\end{theorem}
\begin{proof}
$H$ is positive definite because its diagonal term is positive definite and
its exposure term is positive semidefinite. The continuous strictly convex
objective on the compact convex box therefore has a unique minimizer.
A point minimizes a differentiable convex function on the box exactly when
$\langle Hx-\lambda,y-x\rangle\geq0$ for every feasible $y$:
necessity follows from feasible one-sided derivatives, and sufficiency from
the supporting-hyperplane inequality. Varying one free coordinate proves the
listed signs, which in turn imply this variational inequality for every $y$.
Substitution of $x=d$ gives the full-bail-in test. Independent authority choices
in all states prove uniqueness of the entire credible rule.

For the stability bound let $x=x(\lambda)$ and
$\widetilde x=x(\widetilde\lambda)$. Apply the two variational inequalities
with each minimizer as the other's comparison point and add them. With
$\Delta=x-\widetilde x$ this gives
\[
 m\|\Delta\|_2^2\leq\Delta^\top H\Delta
          \leq\langle\lambda-\widetilde\lambda,\Delta\rangle
          \leq\|\lambda-\widetilde\lambda\|_2\|\Delta\|_2.
\]
Division when $\Delta\ne0$ proves the assertion; otherwise it is immediate.
\end{proof}
For a specified set $F$ of interior coordinates and set $U$ at their upper
bounds, the remaining coordinates being zero, the candidate is explicitly
$x_F=H_{FF}^{-1}(\lambda_F-H_{FU}d_U)$. Checking its bounds and the displayed
gradient signs certifies it. This describes a finite statewise computation;
the number of failed-bank subsets is exponential in $n$.

\section*{Risk feedback and a quantitative equilibrium certificate}
Normalize each bank's designated debt face to one and assume $d_i^S\leq1$.
Banks choose publicly observable $q_i\in[0,1]$ simultaneously before debt
issuance. Conditional on these choices, bank failures are independent
Bernoulli events. State losses remain the credible $x_i^S$ just constructed.
The loss conditional on bank $i$ failing is
\[
 L_i(q_{-i})=\sum_{B\subseteq N\setminus\{i\}}
 \left(\prod_{j\in B}q_j\prod_{j\in N\setminus(B\cup\{i\})}(1-q_j)\right)
 x_i^{B\cup\{i\}}.
\]
Competitive risk-neutral creditors price its debt at
$Q_i=1-q_iL_i(q_{-i})$. As in the one-bank benchmark, the net private objective
is stipulated to be
\[
 b_iq_i-\tfrac{k_i}{2}q_i^2-q_iL_i(q_{-i}),\qquad b_i>0,\quad k_i>0.
\]
This assumption isolates the funding-discipline channel. It does not derive
private surplus or observable risk choice from a hidden-action contract.

\begin{theorem}[Existence, uniqueness under weak interaction, and iteration]
The risk game has a Nash equilibrium. Its best-response map is
\[
 B_i(q)=\operatorname{proj}_{[0,1]}
             \bigl((b_i-L_i(q_{-i}))/k_i\bigr).
\]
For $i\ne j$ define
\[
 \tau_{ij}=\max_{D\subseteq N\setminus\{i,j\}}
 \left|x_i^{D\cup\{i,j\}}-x_i^{D\cup\{i\}}\right|,
 \qquad r=\max_i\frac{\sum_{j\ne i}\tau_{ij}}{k_i}.
\]
If $r<1$, the equilibrium $q^*$ is unique. Simultaneous best responses
$q^{t+1}=B(q^t)$ from any point of the cube satisfy
\[
 \|q^t-q^*\|_\infty\leq r^t\|q^0-q^*\|_\infty,
 \qquad
 \|q-q^*\|_\infty\leq\frac{\|q-B(q)\|_\infty}{1-r}.
\]
If $r\geq1$, only the existence conclusion is asserted.
\end{theorem}
\begin{proof}
Strict concavity in the bank's own choice gives the projected first-order
condition. Each $L_i$ is a multilinear polynomial, so $B$ is a continuous
self-map of the compact convex cube. Brouwer's fixed-point theorem gives a
fixed point, which is a Nash equilibrium by the best-response formula.

Holding other coordinates fixed, the derivative of $L_i$ with respect to
$q_j$ is an average, over the other banks' independent failure subsets $D$,
of $x_i^{D\cup\{i,j\}}-x_i^{D\cup\{i\}}$. Its absolute value is at most
$\tau_{ij}$, including by continuous extension at boundary choices.
Change coordinates one at a time to obtain
\[
 |L_i(q_{-i})-L_i(\widetilde q_{-i})|
     \leq\sum_{j\ne i}\tau_{ij}|q_j-\widetilde q_j|.
\]
Projection onto $[0,1]$ is one-Lipschitz, so $B$ is $r$-Lipschitz for the
sup norm. If $r<1$, any two fixed points coincide. Successive iterates have
geometrically shrinking differences and are Cauchy; their limit is a fixed
point by continuity. Comparing an iterate to that fixed point proves the
first bound by induction. Finally
$\|q-q^*\|\leq\|q-B(q)\|+r\|q-q^*\|$ proves the residual certificate.
\end{proof}

\begin{example}[Two banks and exact endogenous state probabilities]
Let both banks have unit shortfalls and caps and unit fiscal weights. In a
singleton failure set take $H=[1]$, yielding $x=1$. In the joint state take
\[
 H=\begin{pmatrix}2&1\\1&2\end{pmatrix}
       =I+\begin{pmatrix}1\\1\end{pmatrix}\begin{pmatrix}1&1\end{pmatrix}.
\]
This is an admissible exposure specification and its credible minimizer is
$(1/3,1/3)$. With $b_1=b_2=6/5$ and $k_1=k_2=2$, one has
\[
 L_i(q_{-i})=1-\tfrac23q_j,\qquad
 B_i(q)=\tfrac1{10}+\tfrac13q_j,\qquad r=\tfrac13.
\]
The response stays strictly inside the cube, and the unique equilibrium is
$q_1=q_2=3/20$. Conditional expected losses are $9/10$ and debt prices are
$1-(3/20)(9/10)=173/200$. Thus the joint failure probability and creditor
losses are determined together, despite statewise resolution primitives
being fixed. Compared with an otherwise identical isolated bank's risk
$1/10$, the expected dilution of bail-in in the joint state raises risk.
\end{example}

\section*{Unresolved part}
Version 2 endogenizes independent failure probabilities and permits specified
cross-creditor exposure costs. It still assumes the quadratic spillover
function, observable pre-issuance risk, fixed claim priorities and caps, and
fixed state primitives. It neither derives a network clearing mechanism nor
models runs, hidden post-issuance risk, endogenous exposures, liquidity or
transfer of critical functions. Within this benchmark the credible rule is
unique, so there is no nontrivial choice among credible rules without adding
instruments that alter primitives. Welfare-optimal credible instrument design
and empirical identification in the broader catalogue remain unresolved.
The single-price nonidentification result does not dispute strategies using
additional claim classes and identifying restrictions.
\section{Completion Estimate}
35\%

This is a heuristic estimate for the full catalogue target. The new all-state
multibank characterization and risk-equilibrium certificate extend the
one-bank result, but leave clearing, runs, hidden incentives and general
welfare-optimal credible design open.

\begin{thebibliography}{9}
\bibitem{bis} A. Di Stefano, Y. Lengwiler and K. Rishabh (2026), The Credibility of Bail-in, BIS Working Paper 1356, official abstract and findings. Primary source cited in version 1: \url{https://www.bis.org/publications/working-paper-1356-credibility-bail-in}. The model and proofs here are not attributed to that empirical paper.
\end{thebibliography}
\end{document}
